\section{Формирование синтетического массива данных}

\AKHworries{АХ: для описания алгоритмов используем  algorithm2e с параметрами linesnumbered,ruled,vlined. При этом внутри псевдокода все на английском. Описание алгоритма технически корректным языком. Третируем русский и английский текст.  }

В рамках работы был сделан генератор синтетических тестовых заказов для задачи паллетной упаковки. Идея в том, чтобы получить набор входных данных, на которых можно одинаково проверять разные алгоритмы укладки, при этом не использовать реальные заказы ритейлеров и не копировать их напрямую. Поэтому все идентификаторы товаров (SKU) генерируются искусственно, а сами размеры коробок берутся не из конкретных заказов, а из заранее заданного “профиля”, который имитирует типичные размеры в реальных данных.  

Под “профилем” понимается набор нескольких классов коробок, которые встречаются чаще всего: маленькие (почти кубические), средние, плоские (маленькая высота), высокие и большие. Для каждого класса задаются диапазоны по длине, ширине и высоте и вероятность, с которой он выбирается. Когда нужно сгенерировать коробку для очередного SKU, сначала случайно выбирается один из классов по этим вероятностям, потом размеры выбираются внутри диапазонов. Для кубоподобных коробок дополнительно делается небольшая коррекция, чтобы длина и ширина были близки друг к другу, потому что в реальности такие коробки часто бывают “почти квадратными”. Вес генерируется согласованно с габаритами: грубо оценивается через объём коробки и условную плотность, плюс добавляется небольшой шум, чтобы значения не получались слишком “идеальными”. Процедура генерации размеров по профилю используется во всех группах тестов и приведена в алгоритме.

\begin{algorithm}[H]
\caption{Sample box dimensions from the profile (mixture of families)}
\SetKwInOut{Input}{Input}
\SetKwInOut{Output}{Output}
\Input{$profile$ -- list of families with weights and ranges}
\Output{$(L,W,H)$ -- box dimensions (mm)}
$f \leftarrow WeightedPick(profile.families)$\;
$L \leftarrow UniformInt(f.L_{min}, f.L_{max})$\;
$W \leftarrow UniformInt(f.W_{min}, f.W_{max})$\;
$H \leftarrow UniformInt(f.H_{min}, f.H_{max})$\;
\If{$f.shape = \text{cube-like}$}{
  $s \leftarrow \mathrm{round}((L+W)/2)$\;
  $L \leftarrow \mathrm{clamp}(s + UniformInt(-\delta,\delta), f.L_{min}, f.L_{max})$\;
  $W \leftarrow \mathrm{clamp}(s + UniformInt(-\delta,\delta), f.W_{min}, f.W_{max})$\;
}
\If{$W > L$}{
  swap$(L,W)$\;
}
\Return $(L,W,H)$\;
\end{algorithm}

Данные заказа сохраняются в CSV в том формате, который используется дальше в работе: одна строка соответствует одному SKU и содержит количество коробок этого SKU (Quantity), размеры одной коробки и дополнительные поля. То есть коробки одного SKU считаются одинаковыми по габаритам, что соответствует типичной ситуации, когда товар одного вида имеет одну упаковку. Кроме CSV дополнительно формируется текстовый файл метаданных, где записано, к какой группе относится тест, какой сценарий использовался, сколько получилось коробок и SKU, а также какие параметры паллеты использовались (например, 1200×800 и ограничение по высоте 2000 мм). Фиксированный seed позволяет получать одинаковые тесты при повторном запуске, что удобно для регрессионных проверок. Общая процедура построения набора заключается в последовательной генерации заказов для каждой группы и сценария с последующей записью результатов в файлы. Общий алгоритм построения тестового набора приведён в алгоритме.
\begin{algorithm}[H]
\caption{Generate the full test dataset (groups 1--4)}
\SetKwInOut{Input}{Input}
\SetKwInOut{Output}{Output}
\Input{$seed$; pallet parameters $pallet$; size profile $profile$; scenario list $scenarios$}
\Output{directory $dataset$ with order CSV files and \texttt{meta.txt} files}
$RNG \leftarrow Init(seed)$\;
$dataset \leftarrow \emptyset$\;
\For{$g \leftarrow 1$ \KwTo $4$}{
  \ForEach{$s \in scenarios[g]$}{
    \For{$t \leftarrow 1$ \KwTo $M_{g,s}$}{
      $order \leftarrow GenerateOrder(g, s, pallet, profile, RNG)$\;
      WriteCSV(order)\;
      WriteMeta(order)\;
      AddToDataset(dataset, order)\;
    }
  }
}
\Return $dataset$\;
\end{algorithm}


Сам набор тестов разделён на четыре группы, чтобы покрыть разные ситуации. В первой группе генерируются “обычные” заказы, похожие по структуре на реальные, но с контролируемым количеством коробок. Формирование одного заказа выполняется по единому шаблону, который различается только правилами выбора количеств, размеров и категорий в зависимости от группы. Общий алгоритм генерации одного заказа приведён в алгоритме.

\begin{algorithm}[htbp]
\caption{Generate a single order (groups 1--2)}
\label{alg:gen-order-12}
\SetKwInOut{Input}{Input}
\SetKwInOut{Output}{Output}
\Input{$s$ -- scenario id; $pallet$; $profile$; $RNG$}
\Output{$order$ -- list of SKU-level lines: $(sku, qty, L, W, H, weight, strength, aisle, caustic)$}

$order \leftarrow \emptyset$\;

\BlankLine
\textbf{Group 1: realistic orders with fixed $N$}\;
$N \leftarrow ScenarioParam(s).TotalBoxes$\;
$K \leftarrow SampleKGivenN(N)$\;
$qty[1..K] \leftarrow SampleQuantitiesSumToN(K, N)$\;
$C \leftarrow SampleCategoryCount(K)$\;
$aisle[1..K] \leftarrow AssignCategories(K, C)$\;

\For{$i \leftarrow 1$ \KwTo $K$}{
  $sku \leftarrow NewSyntheticSKU()$\;
  $(L,W,H) \leftarrow SampleDimsFromProfile(profile)$\;
  $weight \leftarrow EstimateWeight(L,W,H)$\;
  $strength \leftarrow SampleStrength()$\;
  $caustic \leftarrow SampleCaustic()$\;
  $order.\mathrm{add}(sku, qty[i], L, W, H, weight, strength, aisle[i], caustic)$\;
}

\BlankLine
\textbf{Group 2: predefined quantity patterns}\;
$qtyPattern \leftarrow ScenarioPattern(s)$\;
$K \leftarrow |qtyPattern|$\;
$C \leftarrow SampleCategoryCount(K)$\;
$aisle[1..K] \leftarrow AssignCategories(K, C)$\;

\For{$i \leftarrow 1$ \KwTo $K$}{
  $sku \leftarrow NewSyntheticSKU()$\;
  $(L,W,H) \leftarrow SampleDimsFromProfile(profile)$\;
  $weight \leftarrow EstimateWeight(L,W,H)$\;
  $strength \leftarrow SampleStrength()$\;
  $caustic \leftarrow SampleCaustic()$\;
  $order.\mathrm{add}(sku, qtyPattern[i], L, W, H, weight, strength, aisle[i], caustic)$\;
}

\Return $order$\;
\end{algorithm}



\begin{algorithm}[htbp]
\caption{Generate a single order (groups 3--4)}
\label{alg:gen-order-34}
\SetKwInOut{Input}{Input}
\SetKwInOut{Output}{Output}
\Input{$s$ -- scenario id; $pallet$; $profile$; $RNG$}
\Output{$order$ -- list of SKU-level lines: $(sku, qty, L, W, H, weight, strength, aisle, caustic)$}

$order \leftarrow \emptyset$\;

\BlankLine
\textbf{Group 3: geometry stress tests}\;
$N \leftarrow ScenarioParam(s).TotalBoxes$\;
$K \leftarrow SampleKGivenN(N)$\;
$qty[1..K] \leftarrow SampleQuantitiesSumToN(K, N)$\;
$C \leftarrow SampleCategoryCount(K)$\;
$aisle[1..K] \leftarrow AssignCategories(K, C)$\;

\For{$i \leftarrow 1$ \KwTo $K$}{
  $sku \leftarrow NewSyntheticSKU()$\;
  $(L,W,H) \leftarrow SampleDimsFromGeometryScenario(s)$\;
  $weight \leftarrow EstimateWeight(L,W,H)$\;
  $strength \leftarrow SampleStrength()$\;
  $caustic \leftarrow SampleCaustic()$\;
  $order.\mathrm{add}(sku, qty[i], L, W, H, weight, strength, aisle[i], caustic)$\;
}

\BlankLine
\textbf{Group 4: multi-pallet orders with categories}\;
$P \leftarrow ScenarioParam(s).TargetPallets$\;
$C \leftarrow ScenarioParam(s).CategoryCount$\;
$shares[1..C] \leftarrow SampleCategoryShares(s, C)$\;
$V \leftarrow ComputeTargetVolume(P, pallet)$\;

\For{$c \leftarrow 1$ \KwTo $C$}{
  $V_c \leftarrow shares[c] \cdot V$\;
  $catVol \leftarrow 0$\;

  \While{$catVol < V_c$}{
    $sku \leftarrow NewSyntheticSKU()$\;
    $(L,W,H) \leftarrow SampleDimsFromProfile(profile)$\;
    $qty \leftarrow SampleQtyForCategory(shares[c])$\;
    $weight \leftarrow EstimateWeight(L,W,H)$\;
    $strength \leftarrow SampleStrength()$\;
    $caustic \leftarrow SampleCaustic()$\;
    $order.\mathrm{add}(sku, qty, L, W, H, weight, strength, c, caustic)$\;
    $catVol \leftarrow catVol + qty \cdot Volume(L,W,H)$\;
  }
}

\Return $order$\;
\end{algorithm}

Для этого задаётся общее число коробок N (10, 20, …, 100), а затем выбирается число SKU K, которое обычно увеличивается вместе с N, но с небольшими случайными отклонениями. После этого N коробок распределяются между K SKU так, чтобы у большинства SKU было 1–2 коробки, а у небольшой части SKU встречалось заметно больше коробок. Это распределение строится так, чтобы сохранялась фиксированная сумма $\sum_i Quantity_i = N$ и при этом присутствовал небольшой хвост крупных значений. Соответствующая процедура приведена в алгоритме.

\begin{algorithm}[H]
\caption{Distribute box counts across SKUs with a fixed total}
\SetKwInOut{Input}{Input}
\SetKwInOut{Output}{Output}
\Input{$K$ -- number of SKUs; $N$ -- total number of boxes}
\Output{$qty[1..K]$ such that $\sum_i qty[i] = N$}
\For{$i \leftarrow 1$ \KwTo $K$}{
  $qty[i] \leftarrow 1$\;
}
$remaining \leftarrow N - K$\;
\If{$remaining \le 0$}{
  \Return $qty$\;
}
Shuffle indices $1..K$ into $perm[1..K]$\;
\For{$t \leftarrow 1$ \KwTo $K$}{
  $weights[perm[t]] \leftarrow 1/(t^\alpha)$\;
}
\For{$r \leftarrow 1$ \KwTo $remaining$}{
  $i \leftarrow WeightedRandomIndex(weights)$\;
  $qty[i] \leftarrow qty[i] + 1$\;
}
\Return $qty$\;
\end{algorithm}
Такой “хвост” нужен, потому что в реальности иногда бывают позиции, которые берут пачками. Эта группа нужна в первую очередь, чтобы посмотреть, как алгоритмы масштабируются по N, и насколько стабильно они работают на “обычном” входе.


Вторая группа предназначена для особых случаев, где распределение количества коробок по SKU задаётся заранее. Например, заказ из одного товара с большим количеством коробок, несколько SKU с одинаково большим количеством, или наоборот заказ из большого числа уникальных SKU, где почти везде по одной коробке. Также добавляется смешанный случай, когда есть два SKU с большим количеством коробок и длинный “хвост” из множества SKU по 1–2 коробки. Эти тесты полезны тем, что многие эвристики начинают вести себя плохо именно на таких структурах: либо резко ухудшается качество решения, либо растёт время работы, либо появляются ошибки в обработке данных.

В третьей группе основное внимание уделяется геометрии коробок. Там структура заказа по количеству коробок может быть похожей на группу 1, но сами размеры генерируются специальными режимами: все коробки большие, все коробки маленькие, смесь больших и маленьких, широкий спектр размеров или смесь форм (узкие длинные коробки, плоские, высокие и т.д.). Такая группа нужна для проверки того, насколько алгоритм устойчив к “неудобным” размерам, когда коробки либо слишком мелкие и их очень много, либо наоборот крупные и быстро заполняют основание паллеты, либо имеют сложное соотношение сторон.

Четвёртая группа сделана для многопаллетных заказов и учитывает категории товаров (aisle/департамент). Здесь цель тестов двойная: с одной стороны алгоритм должен упаковать всё в минимальное число паллет при ограничении по высоте, а с другой стороны желательно, чтобы товары одной категории не были размазаны по нескольким паллетам. Чтобы получать заказы примерно масштаба 2–10 паллет, используется генерация по целевому объёму: задаётся целевой уровень “сколько паллет должно получиться примерно”, переводится в целевой суммарный объём груза (с коэффициентом заполнения), а затем объём распределяется между категориями по заданному паттерну. Паттерны могут быть разными: одна большая категория и много маленьких, либо категории примерно равны, либо несколько средних и длинный хвост малых. После этого внутри каждой категории создаются SKU и количества коробок до тех пор, пока суммарный объём этой категории не приблизится к своей доле. В итоге получается заказ, где категории представлены осмысленно и их действительно важно “держать вместе”, что позволяет потом сравнивать алгоритмы не только по числу паллет, но и по качеству распределения категорий.

Таким образом, генератор даёт библиотеку тестов, которая одновременно покрывает “обычные” заказы, крайние паттерны по количествам, сложную геометрию и многопаллетные случаи с категориями. Это позволяет использовать один и тот же набор входных данных для сравнения разных алгоритмов упаковки по времени работы, качеству решения и устойчивости к нестандартным ситуациям.

Чтобы проверить, насколько генератор подходит для массового прогона тестов, была проведена оценка времени генерации для каждой группы. Измерение выполнялось с помощью Stopwatch в C\# и включало полный цикл генерации заказов внутри группы. Результаты показали, что генерация выполняется быстро: для группы 1 было сгенерировано 50 заказов (2750 коробок) за 147 мс, для группы 2 — 24 заказа (1684 коробки) за 61 мс, для группы 3 — 45 заказов (2850 коробок) за 110 мс, для группы 4 — 18 заказов (5201 коробка) за 88 мс. Таким образом, генерация данных занимает доли секунды и не является узким местом при экспериментальном тестировании алгоритмов упаковки.

\begin{table}[h]
\centering
\caption{Performance of the dataset generator for different test groups}
\label{tab:gen_perf}
\begin{tabular}{|c|r|r|r|}
\hline
\textbf{Group} & \textbf{Orders} & \textbf{Boxes} & \textbf{Total time (ms)} \\
\hline
1 & 50 & 2750 & 147 \\
\hline
2 & 24 & 1684 & 61 \\
\hline
3 & 45 & 2850 & 110 \\
\hline
4 & 18 & 5201 & 88 \\
\hline
\end{tabular}
\end{table}

\AKHworries{АХ:  алгоритм Андрея .  }